%% kapa: exam=1 round=36 year=2025 period=1 form=A subject=economics subject_ko=경제학원론 seq=39 number=79
%% source: Q-Net 감정평가사 기출문제 (cst00309) artlSeq=5231525 file=raw/1008/5231525/2025년 제36회 감정평가사 1차시험 1교시 문제지 원본.pdf
%% conversion: kordoc 4.24.51 -> tools/gen1.py (auto); review=proofread-sonnet; flags=none
\begin{question}{79}
아래의 거시경제모형 초기 균형에서 외생적 정부지출( $\alpha$ )을 1만큼 증가시킬 때, 거시경제 변수의 변화에 관한 설명으로 옳지 않은 것은? (단, $Y$ 는 국민소득, $C$ 는 소비, $I$ 는 투자, $G$ 는 정부지출, $T$ 는 조세)
\begin{framedtext}
$\circ$ $Y=C+I+G$ \par
$\circ$ $C=1+0.5(Y-T)$ \par
$\circ$ $I=200-T$ \par
$\circ$ $T=0.5Y+\alpha$ \par
$\circ$ $G=0.5Y+\alpha$
\end{framedtext}
\begin{choices}
  \item 국민소득은 증가한다.
  \item 투자가 감소한다.
  \item 소비는 감소한다.
  \item 균형재정을 달성한다.
  \item 정부지출은 증가한다.
\end{choices}
\end{question}
\answer{1}
